\section*{7. Methodik}
Nach einer ausführlichen Literaturrecherche folgt die Analyse der Daten und Prozesse, welche beim Monitoring Anwendung finden. Anschließend werden die Anforderungen der bisherigen Lösung und die Erwartungen an die neue Lösung analysiert und festgehalten, um sie zur Evaluation nutzen zu können.

Das Dashboard muss konzipiert und entwickelt werden. Dafür werden verschiedene Tools bewertet und eine Skizze mit den wichtigsten Inhalten und Vorgehensweisen aufgestellt. Daraufhin können das Dashboard anhand der Skizze aufgebaut und die entsprechenden Daten eingepflegt werden.

Die Validierung erfolgt mittels Interviews mit den Teammitgliedern, welche die Lösung verwenden werden. Zur Erstellung des Fragenkatalogs für die Interviews ist eine weitere Literaturrecherche notwendig. Hierfür müssen mögliche Antworten und Ergebnisse vorher in Betracht gezogen werden, um die Fragen auf gewünschte Metriken zu reduzieren. Nach Erstellung des Fragenkatalogs und Fertigstellung des Dashboards werden die Interviews durchgeführt und die Ergebnisse ausgewertet.